\section{Taxonomy of XRL Methods}
\label{sec:taxonomy_xrl_methods}

Now that the concept of \gls{xai} has been defined, we turn to its application in \gls{rl}, known as \gls{xrl}.
As in \gls{xai}, explanations in \gls{xrl} are produced by explainability methods and can be described in terms of their subject, source, and form.
This section reviews the possible variations of these three aspects in the specific context of \gls{rl}.

\subsection{Subject}
\label{subsec:subject}

The first question to address is the subject of explanation in the context of \gls{rl}.
To answer this question, it is necessary to identify the different components involved in a \gls{rl} system and determine which of them constitutes the object of explanation.
A \gls{rl} system can be described through two main components: the environment and the policy.

The environment defines the dynamics in which the agent evolves, including the possible states, observations, actions, and transitions.
Environments are simulations designed by humans, and their dynamics are explicitly defined.
Therefore, it is assumed that the environment is known to the observer and does not require explanation.

The main object of interest is therefore the agent's policy.
The policy represents the decision-making mechanism of the agent.
It is a function that maps observations to actions and is learned through interaction with the environment in order to maximize the expected cumulative reward.
The policy is the component that encodes the learning process and that \gls{xrl} aims to explain.

Explaining a policy introduces additional challenges compared with classical \gls{xai} settings.
In traditional \gls{xai} problems, the objective is to explain an isolated decision, such as determining which features of an image led to a classification result.
In contrast, \gls{rl} introduces a temporal dimension into the decision-making process.
An action is not only determined by the current observation but also influences future states and rewards.
Therefore, understanding the behavior of an agent may require considering a sequence of interactions between the agent and its environment.

This temporal dimension leads to a distinction between two objectives of explanation in \gls{xrl}:
\begin{itemize}
    \item \textbf{Decision explanations:} focus on identifying the elements of the current observation that influenced the action selected by the policy.
    This type of explanation is close to classical \gls{xai} approaches, as it analyzes the relationship between inputs and outputs while abstracting away the sequential nature of the decision process.

    \item \textbf{Strategy explanations:} aim to characterize the agent's behavior over time.
    They focus on sequences of decisions and interactions with the environment in order to understand how the agent achieves its long-term objective.
    This type of explanation is specific to \gls{rl}, as it addresses the temporal and goal-oriented nature of the learning process.
\end{itemize}

These two objectives should not be considered as strictly separated categories, but rather as two extremes of a continuum.
Indeed, understanding the elements of an observation that influence a specific decision can help the observer generalize this information to infer broader aspects of the agent's behavior.
Conversely, understanding the agent's overall strategy can provide insights into the reasoning behind individual decisions.

\subsection{Source}
\label{subsec:source}

We identify three sources of explanation in \gls{rl}, all obtained during or after the learning process and containing information about the agent:
\begin{itemize}
    \item \textbf{Policies:} policies are functions that generate a distribution over the action space for a given observation.
    The objective of \gls{rl} is to learn a policy that maximizes the expected cumulative future rewards.
    During the learning phase, the policy is iteratively improved through interactions with the environment.

    \item \textbf{Trajectories:} trajectories represent sequences of successive interactions between the agent and the environment.
    At each step, the agent selects an action based on its observation, which modifies the environment and produces a reward as well as a new observation.
    This process continues until the end of the trajectory.

    \item \textbf{Value functions:} value functions provide predictions about the future outcome of a trajectory from a given observation.
    The most commonly used value function is the critic function, which estimates the expected sum of future rewards.
    This prediction guides the learning process, as the agent updates its policy to maximize the estimated future return.
    However, other types of value functions can also be used as sources of explainability.
\end{itemize}

\subsection{Form}
\label{subsec:form}

We identify six forms of explanation in the \gls{rl} setting:
\begin{itemize}
    \item \textbf{Policy Model:} an interpretable model of the agent's policy function.
    
    \item \textbf{Saliency Map:} a weighting of the input features representing their influence on the output of the policy function.
    
    \item \textbf{Counterfactual:} a modified observation obtained by perturbing the original observation in order to produce a different action from the agent.
    
    \item \textbf{Agent Trajectory:} a sequence of successive interactions between the agent and the environment.
    
    \item \textbf{Trajectory Prediction:} a prediction about the future evolution of the trajectory produced by a value function.
    
    \item \textbf{Interaction Model:} an interpretable model of the interactions between the agent and the environment.
\end{itemize}

Among the three components characterizing an explanation, the form is the one that varies most significantly across explainability methods.
While subjects and sources tend to be relatively stable in the \gls{xrl} literature, the diversity in how explanations are presented is considerably greater.
The list above is by no means definitive and may evolve as new approaches emerge.
For this reason, the following section is organized according to the form of explanations and presents existing explainability methods based on this criterion.
